% === DOCBLOCK 01 BEGIN: Core packages ===
\usepackage[letterpaper,left=1in,right=1in,top=1in,bottom=1.8in,footskip=0.8in,headheight=24pt,headsep=0pt]{geometry} % page size + margins; geometry is your “single source of truth”
\usepackage{setspace}  												% \singlespacing etc.; controls baseline spacing cleanly
\usepackage{eso-pic}   												% place objects on page background (headers/footers overlays, etc.)
\usepackage{enumitem}  												% precise list control (indent, label, spacing, nesting depth)
\setlistdepth{6}       												% allow deeper list nesting (enumitem must be loaded first)
\usepackage{indentfirst}		 									% indent first paragraph after headings (LaTeX default is no indent)
% === DOCBLOCK 01 END ===

% === DOCBLOCK 02 BEGIN: Fonts + bullets + graphics + underlining + tables ===

% --- Fonts (XeLaTeX) ---------------------------------------------------------
\usepackage{fontspec}  												% use system fonts (Times New Roman, Arial, etc.) with XeLaTeX
\defaultfontfeatures{Ligatures=TeX} 								% traditional TeX ligature mapping (`` '' --- etc.)
\setmainfont{Times New Roman} 										% main (roman) text font
\setsansfont{Arial}       					    					% sans-serif font
\setmonofont{Courier New}     										% monospace / code font
% --- CJK (Chinese) -----------------------------------------------------------
\usepackage{xeCJK}													% Use xeCJK so XeLaTeX applies correct CJK spacing/line-breaking rules.
\setCJKmainfont{Noto Serif SC}
\newcommand{\DocDefaultFont}{\normalfont\normalsize\selectfont} 	% handy “reset to default font” macro

% --- Graphics + underline ----------------------------------------------------
\usepackage{graphicx}         										% \includegraphics
\usepackage[normalem]{ulem}   										% \uline underline; [normalem] prevents \emph from becoming underline
\graphicspath{{media/}}       										% search <root>/media for graphics, so you can write \includegraphics{image4}

% --- Doc-controlled lengths --------------------------------------------------
\newlength{\DocDefaultBaselineSkip} 								% custom baseline skip length you control elsewhere
\newlength{\DocListOuterSep}        								% custom list outer separation length you control elsewhere

% --- Tables (pandoc/booktabs style) -----------------------------------------
\usepackage{booktabs}         										% \toprule \midrule \bottomrule for professional tables
\usepackage{array}            										% column modifiers like >{...}, plus \arraybackslash
\usepackage{calc}             										% length arithmetic like (\columnwidth - 2\tabcolsep) * ...
\newcommand{\real}[1]{#1}     										% helper for pandoc-generated decimal factors in calc expressions
% === DOCBLOCK 02 END ===

% === DOCBLOCK 03 BEGIN: List policy ===
\newcommand{\ApplyDocListSpacing}{%
	% Make all itemize lists single-spaced (while body stays double-spaced)
	\setlist[itemize]{%
		before=\begin{singlespace*},
			after=\end{singlespace*},
		itemsep=0pt, parsep=0pt, topsep=\baselineskip,	partopsep=0pt,
	}
	\setlist[itemize,1]{label=\textbullet, left=0.5in..0.75in}
	\setlist[itemize,2]{label=\textbullet, left=0in..0.25in}
	\setlist[itemize,3]{label=\textbullet, left=0in..0.25in}
	\setlist[itemize,4]{label=\textbullet, left=0in..0.25in}
	% Global enumerate policy
	\setlist[enumerate]{before=\begin{singlespace*}, after=\end{singlespace*}, itemsep=0pt, parsep=0pt, topsep=0pt, partopsep=0pt, align=right, labelsep=0.12in }
	% Level 1 (outer list)
	\setlist[enumerate,1]{topsep=\baselineskip, labelindent=0pt, labelwidth=0.25in, leftmargin=0.75in, label=\arabic*.}
	% Levels 2–4 (nested)
	\setlist[enumerate,2]{labelindent=0pt, labelwidth=0.13in, leftmargin=0.25in, label=(\alph*), topsep=\baselineskip, after=\vspace{-\baselineskip}\end{singlespace*}}
	\setlist[enumerate,3]{labelindent=0pt, labelwidth=0.13in, leftmargin=0.25in, label=\roman*., topsep=\baselineskip, after=\vspace{-\baselineskip}\end{singlespace*}}
	\setlist[enumerate,4]{labelindent=0pt, labelwidth=0.13in, leftmargin=0.25in, label=\Alph*., topsep=\baselineskip, after=\vspace{-\baselineskip}\end{singlespace*}}
}
% === DOCBLOCK 03 END ===

% === DOCBLOCK 04 BEGIN: Page numbers + headings ===
\newsavebox{\PageNumBox}
\newcommand{\EnableBottomCenterPageNumbers}{%
	\AddToShipoutPictureFG{%
		\AtPageLowerLeft{%
			\begingroup
			\sbox{\PageNumBox}{\rmfamily\normalsize\thepage}%
			% bottom of numeral is exactly 1in above paper bottom:
			\raisebox{\dimexpr 1in+\dp\PageNumBox\relax}{%
				\makebox[\paperwidth][c]{\usebox{\PageNumBox}}%
			}%
			\endgroup
		}%
	}%
}

% Optional (handy later)
\newcommand{\DisableBottomCenterPageNumbers}{\ClearShipoutPictureFG}
% Don’t let class/page styles print anything (we’re doing it ourselves).
\pagestyle{empty}
\newcommand{\DocTitle}{Gesture-Efficient Global Language Input Methods}
\newcommand{\DocAuthor}{Daniel Mailman}
\newsavebox{\TitleBox}
\savegeometry{NormalPageGeom}
\newgeometry{left=1in,right=1in,top=0in,bottom=1.8in,footskip=0.8in,headheight=24pt,headsep=0pt}
\savegeometry{TopEdgePageGeom}
\loadgeometry{NormalPageGeom}
\newcommand{\UseNormalPageGeom}{\loadgeometry{NormalPageGeom}}
\newcommand{\UseTopEdgePageGeom}{\loadgeometry{TopEdgePageGeom}}
\newcommand{\VSpaceFromTop}[1]{\vspace*{\dimexpr #1-\topskip\relax}} 	% Place baseline at an absolute distance from physical top (TopEdgePageGeom assumed).
% --- Front-matter heading: top of CAPS exactly 2in from physical top ----------
\newlength{\FrontCapHeight}
\settoheight{\FrontCapHeight}{\normalsize X}
\newcommand{\FrontHeadingTwoInches}[1]{%
	\begingroup
	\singlespacing
	\sbox{\TitleBox}{\normalsize\MakeUppercase{#1}}%
	% NormalPageGeom: top of text area is 1in from physical top.
	% Want TOP OF CAPS at 2in => baseline at (2in - capheight).
	% Relative to text-area top: (2in - capheight) - 1in = 1in - capheight.
	\vspace*{\dimexpr 1in - \FrontCapHeight - \topskip\relax}%
	\noindent\makebox[\linewidth][c]{\usebox{\TitleBox}}\par
	\endgroup
}
% === DOCBLOCK 04 END ===

% === DOCBLOCK 05 BEGIN: Chapter headings ===
% ============================================================================
% Main-matter chapter headings (book.cls)
%   - Header at 2" from top (reuse \FrontHeadingTwoInches)
%   - One single-spaced blank line between header and title
%   - Two single-spaced blank lines between title and body
%   - Header + title in ALL CAPS, default font/size
%   - Uses \@chapapp so appendices become "APPENDIX A" automatically
% ============================================================================
\makeatletter

\def\@makechapterhead#1{%
	\begingroup
	\singlespacing
	% (2) header position + (5) caps in default font
	\FrontHeadingTwoInches{\@chapapp\space\thechapter}%
	% (3) one single-spaced blank line
	\vspace{\baselineskip}%
	% (5) title all caps, default font/size, centered
	\noindent\makebox[\linewidth][c]{\normalsize\normalfont\MakeUppercase{#1}}\par
	% (4) two single-spaced blank lines before body
	\vspace{2\baselineskip}%
	\endgroup
	% restore doc spacing for body text
	\DocDefaultSpacing
}

\def\@makeschapterhead#1{%
	\begingroup
	\singlespacing
	\FrontHeadingTwoInches{#1}%
	\vspace{2\baselineskip}%
	\endgroup
	\DocDefaultSpacing
}

\makeatother
% === DOCBLOCK 05 END ===

% === DOCBLOCK 06 BEGIN: Default spacing + helpers ===
\newlength{\SavedFootskip}
\newcommand{\RestoreFootskipAtShipout}{\setlength{\footskip}{\SavedFootskip}}

\newcommand{\DocDefaultSpacing}{%
	\setstretch{2}\selectfont
	\setlength{\parindent}{0.5in}%
	\setlength{\parskip}{0pt}%
	\setlength{\DocDefaultBaselineSkip}{\baselineskip}%
	\setlength{\DocListOuterSep}{0.5\DocDefaultBaselineSkip}%
	\setlength{\DocTOCBlankSkip}{0.5\DocDefaultBaselineSkip}%
}


\newcommand{\AbsDedAckPage}[2]{%
	% --- Abstract / Dedication / Acknowledgements page ----------------------------
	% #1 = title text, #2 = content (usually \input{...})
	\clearpage
	\UseNormalPageGeom
	\thispagestyle{empty}% you handle page numbers via your shipout overlay
	\addcontentsline{toc}{frontmatter}{#1}%
	\FrontHeadingTwoInches{#1}%
	\vspace{2\baselineskip}% exactly two blank lines after the heading
	 \DocDefaultSpacing   % <-- guarantee body text returns to default here
	#2%
}
% === DOCBLOCK 06 END ===
\newcommand{\VSpaceDoubleLines}[1]{\begingroup\begin{spacing}{2}\vspace{#1\baselineskip}\end{spacing}\endgroup}
\newlength{\TitleTopWhiteSpace}
\setlength{\TitleTopWhiteSpace}{2in}
\newlength{\TitleTopAdjust}
\setlength{\TitleTopAdjust}{-13pt} % your desired value

% === DOCBLOCK 07 BEGIN: Front matter 01 committee ===
%% ============================================================================
%% Front Matter 01 (1 Page) - Committee and University Representative Page
%% ============================================================================
\newcommand{\CommitteeAndRepPage}{%
	\DisableBottomCenterPageNumbers
	\newcommand{\CommitteeLeftOneName}{Dalei Wu}
	\newcommand{\CommitteeLeftOneTitle}{Associate Professor of Computer Science}
	\newcommand{\CommitteeLeftOneRole}{(Chair)}
	\newcommand{\CommitteeRightOneName}{Cuilan Gao}
	\newcommand{\CommitteeRightOneTitle}{Associate Professor of Mathematics}
	\newcommand{\CommitteeRightOneRole}{(Committee Member)}
	\newcommand{\CommitteeLeftTwoName}{Yingfeng Wang}
	\newcommand{\CommitteeLeftTwoTitle}{Associate Professor of Computer Science}
	\newcommand{\CommitteeLeftTwoRole}{(Committee Member)}
	\newcommand{\CommitteeRightTwoName}{Yu Liang}
	\newcommand{\CommitteeRightTwoTitle}{Professor of Computer Science}
	\newcommand{\CommitteeRightTwoRole}{(Committee Member)}
	\newlength{\CommitteeLeftW}
	\newlength{\CommitteeRightW}
	\newlength{\CommitteeGap}
	\newcommand{\SetupCommitteeWidths}{%
		\settowidth{\CommitteeLeftW}{\normalsize\CommitteeLeftOneTitle}%
		\settowidth{\CommitteeRightW}{\normalsize\CommitteeRightOneTitle}%
		\setlength{\CommitteeGap}{\dimexpr\linewidth-\CommitteeLeftW-\CommitteeRightW\relax}%
	}
	\clearpage
	\UseTopEdgePageGeom
	\thispagestyle{empty}
	\SetupCommitteeWidths
	\begingroup
	\sbox{\TitleBox}{\normalsize\MakeUppercase{\DocTitle}}%
	\VSpaceFromTop{\TitleTopWhiteSpace + \TitleTopAdjust + \ht\TitleBox}%
	\noindent\makebox[\linewidth][c]{\usebox{\TitleBox}}\par
	\endgroup
	\VSpaceDoubleLines{4}
	{\centering
		By\\[\baselineskip]
		\DocAuthor\par
	}
	\VSpaceDoubleLines{4}
	\noindent
	\begin{minipage}[t]{\CommitteeLeftW}\raggedright
		\singlespacing
			\CommitteeLeftOneName\\
			\CommitteeLeftOneTitle\\
			\CommitteeLeftOneRole
	\end{minipage}%
	\hspace{\CommitteeGap}%
	\begin{minipage}[t]{\CommitteeRightW}\raggedright
		\singlespacing
			\CommitteeRightOneName\\
			\CommitteeRightOneTitle\\
			\CommitteeRightOneRole
	\end{minipage}\par
	\vspace{1\baselineskip}
	\noindent
	\begin{minipage}[t]{\CommitteeLeftW}\raggedright
		\singlespacing
			\CommitteeLeftTwoName\\
			\CommitteeLeftTwoTitle\\
			\CommitteeLeftTwoRole
	\end{minipage}%
	\hspace{\CommitteeGap}%
	\begin{minipage}[t]{\CommitteeRightW}\raggedright
		\singlespacing
			\CommitteeRightTwoName\\
			\CommitteeRightTwoTitle\\
			\CommitteeRightTwoRole
	\end{minipage}\par
}%
% === DOCBLOCK 07 END ===

% === DOCBLOCK 08 BEGIN: Front matter 02 title ===
% ============================================================================
% Front Matter 02 (1 Page) - Title Page (number "ii")
% ============================================================================
\newcommand{\UniversityName}{The University of Tennessee at Chattanooga}
\newcommand{\UniversityLocation}{Chattanooga, Tennessee}
\newcommand{\AwardMonthYear}{May 2026} % No comma; allowed: May/August/December
\newcommand{\DegreeName}{Doctor of Philosophy}
\newcommand{\DocumentType}{Dissertation} % "Thesis" or "Dissertation"
\newcommand{\CopyrightYear}{\the\year}
\newcommand{\TitlePage}{%
	\clearpage
	\EnableBottomCenterPageNumbers
	\UseTopEdgePageGeom
	\thispagestyle{empty}

	% ---- Title: baseline positioned from physical top edge (your existing mechanism) ----
	\begingroup
	\singlespacing
	\sbox{\TitleBox}{\normalsize\MakeUppercase{\DocTitle}}%
	\VSpaceFromTop{\TitleTopWhiteSpace + \TitleTopAdjust + \ht\TitleBox}%
	\noindent\makebox[\linewidth][c]{\usebox{\TitleBox}}\par
	\endgroup

	% ---- Line-counted layout below ----
	\begin{spacing}{1}
		{\centering
			% 10 lines between title and By
			\vspace{10\baselineskip}

			By\par
			% 1 line between By and name
			\vspace{1\baselineskip}

			\DocAuthor\par
			% 10 lines between name and submission block
			\vspace{10\baselineskip}

			% Submission block: single spaced
			\begin{spacing}{1}
				A \DocumentType\ Submitted to the Faculty of the University of\\
				Tennessee at Chattanooga in Partial\\
				Fulfillment of the Requirements of the Degree\\
				of \DegreeName\par
			\end{spacing}

			% 4 lines before University/Location
			\vspace{4\baselineskip}

			\UniversityName\par
			\UniversityLocation\par

			% 1 line before date
			\vspace{1\baselineskip}

			\AwardMonthYear\par
		}
	\end{spacing}

	\UseNormalPageGeom
}%
% === DOCBLOCK 08 END ===

% === DOCBLOCK 09 BEGIN: Front matter 03 copyright ===
% ============================================================================
% Front Matter 03 (1 Page) Copyright Page (number "iii")
% ============================================================================
\newcommand{\CopyrightPage}{%
		\clearpage
		% For this page we measure from the physical top edge.
		\UseTopEdgePageGeom
		\thispagestyle{empty}
		\begingroup
		% Avoid first-line/topskip surprises: for this page, treat "top of text area"
		% as the physical top edge.
		\topskip=0pt
		% Build the 3-line block exactly as required.
		% NOTE: line 2 is "By Daniel Mailman" on ONE line (copyright page only).
		\setbox0=\vbox{%
				\hsize=\linewidth
				\centering
				\begin{spacing}{2}
						Copyright \textcopyright\ \CopyrightYear\par
						By \DocAuthor\par
						All Rights Reserved\par
					\end{spacing}
		}%
		% Center the *box* on the physical page:
		% top of box = (paperheight - boxheight)/2
		\vspace*{\dimexpr 0.5\paperheight - 0.5\ht0 - 0.5\dp0\relax}%
		\noindent\box0\par
		\endgroup
}
% === DOCBLOCK 09 END ===

% === DOCBLOCK 10 BEGIN: Glossaries + TOC engine ===
% ============================================================================
% Glossaries (Acronyms) — no-index workflow
% ============================================================================
\usepackage[acronym,nonumberlist,nopostdot]{glossaries-extra}
\makenoidxglossaries 	% Use glossaries' “no-index” mode:

% =============================================================================
% TOC layout: fixed title columns + right-aligned numbers, no overlays
% =============================================================================
\makeatletter

% --- tiny debug helper: writes to .log only (toggle by changing body) ---------
%\newcommand{\TOCDBG}[1]{\typeout{TOCDBG: #1}}
% If you want to turn debug off later, just change to:
\newcommand{\TOCDBG}[1]{}

% --- “Character unit” for column math (measured in the TOC font) --------------
\newlength{\DocTOCChar}          % width of "0"
\newlength{\DocTOCDotW}          % width of '.'

% --- Tunables ----------------------------------------------------------------
\newlength{\DocTOCLeftPad}       % left margin -> start of number-field area
\newlength{\DocTOCTitleGap}      % '.' -> title gap

% --- Derived positions -------------------------------------------------------
\newlength{\DocTOCNumFieldWidth} % width of widest chapter number (e.g. "88")
\newlength{\DocTOCTitlePos}      % where chapter/section titles start
\newlength{\DocTOCSubsecTitlePos}% where subsection titles start

% --- Vertical spacing rule: blank after blocks, not before -------------------
\newlength{\DocTOCBlankSkip}
\newif\ifDocTOCPrevWasSublevel
\DocTOCPrevWasSublevelfalse

\newcommand*\DocTOCMaybeAddBlankSkip{%
	\ifDocTOCPrevWasSublevel
	\addvspace{\DocTOCBlankSkip}%
	\global\DocTOCPrevWasSublevelfalse
	\fi
}

% --- Compute all TOC positions in the actual TOC font ------------------------
\newcommand*\DocInitTOCPositions{%
	\rmfamily\normalfont\normalsize
	% Measure “character” and dot widths in the TOC font
	\setbox0=\hbox{0}\setlength{\DocTOCChar}{\wd0}%
	\setbox0=\hbox{.}\setlength{\DocTOCDotW}{\wd0}%
	% Single-spaced vertical blank
	\setlength{\DocTOCBlankSkip}{\baselineskip}%
	% Tunables: 3 character units
	\setlength{\DocTOCLeftPad}{3\DocTOCChar}%
	\setlength{\DocTOCTitleGap}{3\DocTOCChar}%
	% Number field: allow up to two-digit chapters: "88"
	\setbox0=\hbox{88}\setlength{\DocTOCNumFieldWidth}{\wd0}%
	% Column where chapter and section titles start:
	%   <3spaces> <num field> '.' <3spaces> <Title>
	\setlength{\DocTOCTitlePos}{%
		\dimexpr
		\DocTOCLeftPad
		+ \DocTOCNumFieldWidth
		+ \DocTOCDotW
		+ \DocTOCTitleGap
		\relax}%
	% Subsection titles start 3 spaces further right
	\setlength{\DocTOCSubsecTitlePos}{%
		\dimexpr\DocTOCTitlePos + 3\DocTOCChar\relax}%
	% --- DEBUG readout so we know what we actually computed ---
	\TOCDBG{InitTOC: Char=\the\DocTOCChar, DotW=\the\DocTOCDotW,%
		LeftPad=\the\DocTOCLeftPad, TitleGap=\the\DocTOCTitleGap,%
		NumField=\the\DocTOCNumFieldWidth,%
		TitlePos=\the\DocTOCTitlePos, SubsecPos=\the\DocTOCSubsecTitlePos}%
}

% -----------------------------------------------------------------------------
% CHAPTER entries (includes appendices):
%   "   1.   Title"
%   - Arabic chapter numbers (possibly double digit) are right-aligned in the
%     number field; '.' and title column are fixed.
%   - One single-spaced blank line after each chapter/appendix entry.
% -----------------------------------------------------------------------------
\renewcommand*\l@chapter[2]{%
	\TOCDBG{l@chapter: "#1" page #2; TitlePos=\the\DocTOCTitlePos}%
	\ifnum\c@tocdepth>-1\relax
	\DocTOCMaybeAddBlankSkip
	\begingroup
	\setstretch{1}%
	% Redefine \numberline just for this chapter entry:
	\renewcommand*\numberline[1]{%
		\hb@xt@\DocTOCTitlePos{%
			\hspace*{\DocTOCLeftPad}%         <3 spaces>
			\hb@xt@\DocTOCNumFieldWidth{\hfil ##1}% right-aligned number
			.\hspace*{\DocTOCTitleGap}%       '.' + <3 spaces>
		}%
	}%
	\@dottedtocline{0}{0pt}{\DocTOCTitlePos}{#1}{#2}%
	\addvspace{\DocTOCBlankSkip}% single blank line after chapter
	\global\DocTOCPrevWasSublevelfalse
	\endgroup
	\fi
}

% -----------------------------------------------------------------------------
% SECTION entries:
%   - Unnumbered in display (we hide 1.1 etc.)
%   - Title starts in the SAME column as chapter titles.
%   - No blank lines before/after; only the *next* top-level entry
%     (chapter/divider/frontmatter) will insert a blank after the group.
% -----------------------------------------------------------------------------
\renewcommand*\l@section[2]{%
	\TOCDBG{l@section: "#1" page #2; TitlePos=\the\DocTOCTitlePos}%
	\ifnum\c@tocdepth>0\relax
	\begingroup
	\setstretch{1}%
	% Eat number box but reserve width so text column stays fixed
	\renewcommand*\numberline[1]{\hb@xt@\@tempdima{}}%
	\@dottedtocline{1}{0pt}{\DocTOCTitlePos}{#1}{#2}%
	\global\DocTOCPrevWasSubleveltrue
	\endgroup
	\fi
}

% -----------------------------------------------------------------------------
% SUBSECTION entries:
%   - Unnumbered in display.
%   - Title starts 3 “spaces” to the right of section titles.
% -----------------------------------------------------------------------------
\renewcommand*\l@subsection[2]{%
	\TOCDBG{l@subsection: "#1" page #2; SubsecPos=\the\DocTOCSubsecTitlePos}%
	\ifnum\c@tocdepth>1\relax
	\begingroup
	\setstretch{1}%
	\renewcommand*\numberline[1]{\hb@xt@\@tempdima{}}%
	\@dottedtocline{2}{0pt}{\DocTOCSubsecTitlePos}{#1}{#2}%
	\global\DocTOCPrevWasSubleveltrue
	\endgroup
	\fi
}

% --- Front-matter entry type -------------------------------------------------
\providecommand*{\toclevel@frontmatter}{0}

\providecommand*\l@frontmatter[2]{%
	\TOCDBG{l@frontmatter: "#1" page #2}%
	\DocTOCMaybeAddBlankSkip
	\begingroup
	\setstretch{1}%
	% Front-matter entries: flush-left, no number block.
	\@dottedtocline{0}{0pt}{0pt}{#1}{#2}%
	\endgroup
	\addvspace{\DocTOCBlankSkip}% one blank line after each front-matter entry
}

% --- Divider words like "CHAPTERS", "APPENDICES" -----------------------------
\newcommand*\l@tocdivider[2]{%
	\DocTOCMaybeAddBlankSkip
	\par\noindent\normalfont #1\par
	\addvspace{\DocTOCBlankSkip}%
}

\makeatother

% === DOCBLOCK 10 END ===

% === DOCBLOCK 11 BEGIN: LOF/LOT entry formatting ===
\makeatletter

% Debug helper (can silence later if desired)
\newcommand{\LOXDBG}[1]{\typeout{LOXDBG: #1}}

% Core line formatter for LOF/LOT:
%  - First line flush-left
%  - Wrapped lines hang at 0.5in
%  - Dot leaders to right-aligned page number
\newcommand*\DocLoXLine[2]{%
	\par
	\begingroup
	\parindent=0pt
	\leftskip=0pt
	% continuation lines only:
	\hangindent=0.5in
	\hangafter=1
	% page-number box and dot leaders (LaTeX-native style)
	\rightskip=\@pnumwidth plus 1fil
	\parfillskip=-\@pnumwidth
	\leavevmode
	#1\nobreak
	\leaders\hbox{$\m@th\mkern \@dotsep mu\hbox{.}\mkern \@dotsep mu$}\hfill
	\nobreak\hb@xt@\@pnumwidth{\hfil #2}\par
	\endgroup
}

% LIST OF FIGURES entries
%  - 1 blank line between entries
%  - number and caption inline on first line
%  - hanging indent at 0.5in for wrapped text
\renewcommand*\l@figure[2]{%
	\LOXDBG{figure entry: "#1" page #2}%
	\addvspace{\baselineskip}% one single-spaced blank line between entries
	\begingroup
	% When the .lof line is read, it contains \numberline{<num>}<caption>.
	% Redefine \numberline so it just prints "<num> " inline before the caption.
	\renewcommand\numberline[1]{##1\enspace}%
	\DocLoXLine{#1}{#2}%
	\endgroup
}

% LIST OF TABLES entries (same policy as figures)
\renewcommand*\l@table[2]{%
	\LOXDBG{table entry: "#1" page #2}%
	\addvspace{\baselineskip}% one single-spaced blank line between entries
	\begingroup
	\renewcommand\numberline[1]{##1\enspace}%
	\DocLoXLine{#1}{#2}%
	\endgroup
}

\makeatother
% === DOCBLOCK 11 END ===

% === DOCBLOCK 12 BEGIN: TOC/LOF/LOT pages + symbols ===
\makeatletter
	\newcommand{\AddToCChapterLabel}{%
		\addtocontents{toc}{\protect\setcounter{tocdepth}{2}}%
		\addtocontents{toc}{\protect\contentsline{tocdivider}{CHAPTERS}{}{}}%
	}
	\newcommand{\AddToCAppendixLabel}{%
		\addtocontents{toc}{\protect\setcounter{tocdepth}{0}}%
		\addtocontents{toc}{\protect\contentsline{tocdivider}{APPENDICES}{}{}}%
	}

	\newcommand{\TableOfContentsPage}{%
		\clearpage
		\UseNormalPageGeom
		\thispagestyle{empty}%
		\FrontHeadingTwoInches{TABLE OF CONTENTS}%
		\vspace{2\baselineskip}% <-- two blank lines after the TOC heading
		\begingroup
		\rmfamily\normalfont\normalsize
		\singlespacing
		\setlength{\parskip}{0pt}
		\setlength{\parindent}{0pt}
		\DocInitTOCPositions
		\@starttoc{toc}%
		\DocTOCMaybeAddBlankSkip
		\endgroup
	}
	\newcommand{\ListOfTablesPage}{%
		\clearpage
		\UseNormalPageGeom
		\thispagestyle{empty}%
		\addcontentsline{toc}{frontmatter}{LIST OF TABLES}%
		\FrontHeadingTwoInches{LIST OF TABLES}%
		\begingroup
		\rmfamily\normalfont\normalsize
		\singlespacing
		\setlength{\parindent}{0pt}
		\setlength{\parskip}{0pt}
		\vspace*{\baselineskip}			% 2nd blank line between title and content
		\@starttoc{lot}%
		\endgroup
	}

	\newcommand{\ListOfFiguresPage}{%
		\clearpage
		\UseNormalPageGeom
		\thispagestyle{empty}%
		\addcontentsline{toc}{frontmatter}{LIST OF FIGURES}%
		\FrontHeadingTwoInches{LIST OF FIGURES}%
		\begingroup
		\rmfamily\normalfont\normalsize
		\singlespacing
		\setlength{\parindent}{0pt}
		\setlength{\parskip}{0pt}
		\vspace*{\baselineskip}			% 2nd blank line between title and content
		\@starttoc{lof}%
		\endgroup
	}


	\newcommand{\ListOfSymbolsPage}[1][]{%
		\clearpage
		\UseNormalPageGeom
		\thispagestyle{empty}%
		\addcontentsline{toc}{frontmatter}{LIST OF SYMBOLS}%
		\FrontHeadingTwoInches{LIST OF SYMBOLS}%
		\begin{singlespace}
			\setlength{\parindent}{0pt}
			\printnoidxglossary[type=symbols,style=los-style]
			#1%
		\end{singlespace}
	}
\makeatother
% === DOCBLOCK 12 END ===
% === DOCBLOCK 13 BEGIN: Glossary styles + MainMatterBegin ===
% -ish LOA formatting: "SHORT, Long form" one per line, no page numbers
\newglossarystyle{loa-style}{%
	\setglossarystyle{list}	% Base style: a simple list environment (label + item text)
	\renewcommand*{\glossarysection}[2][]{}% Suppress the default heading normally generated for a glossary.
	\renewcommand*{\glossarypreamble}{}%% Suppress any preamble text inserted before the list.
	\renewcommand*{\glossentry}[2]{%% ##1 = entry label (internal key), ##2 is typically page list (unused here)
		\item[\glsentryshort{##1},] \glsentrylong{##1}%% Print format:%   LABEL:  short form + comma TEXT :  long form
	}%
	\renewcommand*{\glsgroupskip}{}% Suppress extra vertical space between alphabetic groups.
}
% --- Glossaries: symbols ------------------------------------------------------
\newglossary*{symbols}{List of Symbols}
\newglossarystyle{los-style}{%
	\setglossarystyle{list}%
	\renewcommand*{\glossarysection}[2][]{}%
	\renewcommand*{\glossarypreamble}{}%
	\renewcommand*{\glossentry}[2]{%
		\item[\glsentryname{##1},] \glsentrydesc{##1}%
	}%
	\renewcommand*{\glsgroupskip}{}%
}
% Load symbol entries (may be a superset; only referenced ones will print)
% Auto-extracted acronym definitions for UTC dissertation
% Keep this file loaded BEFORE the List of Abbreviations page is printed.

\newacronym{ime}{IME}{Input Method Editor}
\newacronym{ipa}{IPA}{International Phonetic Alphabet}
\newacronym{nlp}{NLP}{Natural Language Processing}
\newacronym{ui}{UI}{User Interface}
\newacronym{utf}{UTF}{Unicode Transformation Format}
\newacronym{qwerty}{QWERTY}{QWERTY keyboard layout}

\newglossaryentry{sym:mu}{
	type=symbols,
	name={\ensuremath{\mu}},
	sort=mu,
	description={Micro (prefix); used in terms such as \emph{\ensuremath{\mu Lex}}}
}
\newglossaryentry{sym:tm}{
	type=symbols,
	name={\textsuperscript{\texttrademark}},
	sort=tm,
	description={Trademark symbol; as in \emph{\ensuremath{\mu Lex}\texttrademark}}
}
\newglossaryentry{sym:approx}{
	type=symbols,
	name={\ensuremath{\approx}},
	sort=approx,
	description={Approximately equal to; used for rough counts}
}

\newglossaryentry{sym:C}{
	type=symbols,
	name={\ensuremath{C}},
	sort=c,
	description={Number of characters entered}
}

\newglossaryentry{sym:G}{
	type=symbols,
	name={\ensuremath{G}},
	sort=g,
	description={Number of gestures performed}
}

\newglossaryentry{sym:CPG}{
	type=symbols,
	name={\ensuremath{C/G}},
	sort=cpg,
	description={Characters per gesture (efficiency metric)}
}

\newglossaryentry{sym:v}{
	type=symbols,
	name={\ensuremath{\mathbf{v}}},
	sort=v,
	description={Vector-valued quantity (e.g., \ensuremath{\mathbf{v}_{\mathrm{lexeme}}})}
}

% --- Symbols used in Chapter 01 (smoke tests) -------------------------------

\newglossaryentry{sym:micro}{%
	type=symbols,
	name={\ensuremath{\mu}},
	sort={mu},
	description={Greek letter mu; used here for the “µ” prefix (e.g., µLex)}
}

\newglossaryentry{sym:lambda}{%
	type=symbols,
	name={\ensuremath{\lambda}},
	sort={lambda},
	description={Greek letter lambda}
}

\newglossaryentry{sym:gpl}{%
	type=symbols,
	name={\ensuremath{\mathrm{GPL}}},
	sort={gpl},
	description={Gestures per Lexeme}
}

\newcommand{\ListOfAbbreviationsPage}[1][]{%
	\clearpage
	\UseNormalPageGeom
	\thispagestyle{empty}%
	% If you haven't defined \l@frontmatter yet, use {chapter} for now (see note below).
	\addcontentsline{toc}{frontmatter}{LIST OF ABBREVIATIONS}%
	\FrontHeadingTwoInches{LIST OF ABBREVIATIONS}%
	\begin{singlespace}
		\setlength{\parindent}{0pt}
		% Primary source: acronyms defined via \newacronym / \newabbreviation.
		% Printed as: "ABBR, Full Term" (no page numbers).
		\printnoidxglossary[type=\acronymtype,style=loa-style]
		% Optional manual additions (normally leave empty):
		#1
	\end{singlespace}
	% Reset "first use" flags so acronyms expand on first use in the main matter.
	\glsresetall
}
% ============================================================================
% Main matter switch (front matter -> chapters)
% ============================================================================
\newcommand{\MainMatterBegin}{%
	\clearpage
	\pagenumbering{arabic}%
	\renewcommand{\thechapter}{\arabic{chapter}}%
	\setcounter{secnumdepth}{3}% number down to subsubsection
	\setcounter{page}{1}%
	\EnableBottomCenterPageNumbers%
	\setcounter{tocdepth}{2}% show sections/subsections in main-matter TOC
}
% === DOCBLOCK 13 END ===
% === DOCBLOCK 14 BEGIN: Section heading formats ===
\makeatletter
\renewcommand\section{%
	\@startsection{section}{1}{\z@}%
	{-3.5ex \@plus -1ex \@minus -.2ex}%
	{2.3ex \@plus .2ex}%
	{\DocDefaultFont\bfseries\raggedright}%
}

\renewcommand\subsection{%
	\@startsection{subsection}{2}{\z@}%
	{-3.25ex \@plus -1ex \@minus -.2ex}%
	{1.5ex \@plus .2ex}%
	{\DocDefaultFont\bfseries\centering}%
}

\renewcommand\subsubsection{%
	\@startsection{subsubsection}{3}{\z@}%
	{-3.25ex \@plus -1ex \@minus -.2ex}%
	{1.5ex \@plus .2ex}%
	{\DocDefaultFont\bfseries\itshape\centering}%
}

\makeatother
% === DOCBLOCK 14 END ===
% === DOCBLOCK 15 BEGIN: Bib + PrintReferences + debug ===
% ============================================================================
% Bibliography (biblatex/biber) — supports "REFERENCES" + "WORKS CONSULTED"
% ============================================================================
\usepackage[backend=biber,style=authoryear,sorting=nyt]{biblatex}
% --- Force DOI/URL onto their own lines in the bibliography ---
\DeclareFieldFormat{doi}{%
	\newline\mkbibacro{DOI}\addcolon\space
	\ifhyperref
	{\href{https://doi.org/#1}{\nolinkurl{#1}}}
	{\nolinkurl{#1}}%
}
\DeclareFieldFormat{url}{%
	\newline\mkbibacro{URL}\addcolon\space\url{#1}%
}
% Point this at your actual .bib file:
% (e.g. references-dissertation.bib, references.bib, etc.)
\addbibresource{REFERENCE.bib}
% Track which entries were actually cited
\DeclareBibliographyCategory{cited}
\AtEveryCitekey{\addtocategory{cited}{\thefield{entrykey}}}
% ============================================================================
% References printing (non-scr, heading 2 inches from top)
% ============================================================================
\newcommand{\PrintReferences}{%
	\nocite{*}%
	% -------------------- Cited works only (REFERENCES) --------------------
	\clearpage
	\UseNormalPageGeom
	\thispagestyle{empty}%
	\addcontentsline{toc}{frontmatter}{REFERENCES}%
	\FrontHeadingTwoInches{REFERENCES}%
	\begingroup
	\begin{singlespace}
		\printbibliography[heading=none,category=cited]
	\end{singlespace}
	\endgroup

	% -------------------- Uncited works only (WORKS CONSULTED) -------------
	\clearpage
	\UseNormalPageGeom
	\thispagestyle{empty}%
	\addcontentsline{toc}{frontmatter}{WORKS CONSULTED}%
	\FrontHeadingTwoInches{WORKS CONSULTED}%
	\begingroup
	\begin{singlespace}
		\printbibliography[heading=none,notcategory=cited]
	\end{singlespace}
	\endgroup
}
\AtBeginDocument{
	\DocDefaultSpacing
	\ApplyDocListSpacing
}
%\makeatletter
%\let\Old@startsection\@startsection
%\def\@startsection#1#2#3#4#5#6{%
%	\typeout{HEADING-DBG (#1): font=\fontname\font}%
%	\Old@startsection{#1}{#2}{#3}{#4}{#5}{#6}%
%}
\makeatother
\makeatletter
\newcommand{\DBGListGlue}{%
	\typeout{LISTDBG depth=\the\@listdepth
		topsep=\the\topsep
		partopsep=\the\partopsep
		itemsep=\the\itemsep
		parsep=\the\parsep
		parskip=\the\parskip
		baselineskip=\the\baselineskip}%
}
\makeatother
%=== DOCBLOCK 15 END ===